\section{Applying a capability predicate to an agent}
A classical state encoder $G:s\mapsto\rho$ and a quantum channel $\Lambda:M_7\to M_7$ have different domains. For example $A(s)A(s)^\dagger/\Tr(A(s)A(s)^\dagger)$ always defines a density matrix when $A(s)\ne0$. Taking $A(s)$ constant makes every input encode the same state. Positivity alone cannot establish that the encoder retains even one task-relevant distinction.

A predictor of a recorded response $q(s)$ can factor through $G$ only if $G(s_1)=G(s_2)$ implies $q(s_1)=q(s_2)$. This is Proposition~\ref{prop:access} applied to the fibres of the encoder. Operational adequacy additionally requires compatible subsequent actions. If a future intervention reveals a difference between two states that $G$ merges, augment the state or restrict the target and domain. A large model can also have a reliable functional capability without its internal variables supporting this particular seven-level representation.

\paragraph{An engineering comparison.}
Record the observation stream and action opportunities at the ordinary task entry point. Vary whether a controlled signal is available, match resources and environmental replies where the design calls for replay, and score on held-out situations. Trace each scored action through the actual policy and its environment reply. A difference between monitored and ablated behaviour can support causal use of that signal under the tested conditions. Reproducing a predicted self-description by following the prompt is a negative control. A proxy-gate crossing alone supplies neither this behavioural evidence nor a phenomenal attribution.

\paragraph{A transfer bound.}
Suppose three reporting schemes assign labels $a(e),b(e),c(e)$ to the same episode $e$. Freeze translations $f:A\to B$, $g:B\to C$ and $h:A\to C$ before testing. Let $M$ count disagreements of $g(f(a(e)))$ with $h(a(e))$. Let $E_f,E_g,E_h$ count translation errors against the corresponding recorded labels on the same list of episodes; in particular $E_g$ compares $g(b(e))$ with $c(e)$. Then
\begin{equation}
 M\le E_f+E_g+E_h .
\label{eq:triangle}
\end{equation}
If all three edges are correct on an episode, the routes both end at its recorded $c(e)$. Summing the contrapositive indicator inequality proves \eqref{eq:triangle} without independence. Eight route disagreements together with $E_f,E_g\le2$ entail $E_h\ge4$. Two routes that always give the same wrong name show why $M=0$ alone does not certify accuracy. The \source{network/PhenomenalAccess.lean}{formal transfer bound} is about paired records. It does not translate experiences between subjects without an independently justified bridge.

\paragraph{Scope of the attribution.}
Experience, a report about experience and a computed label are three distinct targets. A claim about the first needs a stated relation to the two observable targets and a consequence that competing explanations do not also predict. Two-aspect monism is a philosophical interpretation of such a relation. A finite matrix, an accurate decoder, or a trace-preserving map does not entail that an organism or machine has an experience. Normative decisions about welfare require their own premises.
